\hypertarget{walrus-operator-parameter-binding}{%
\chapter{WALRUS OPERATOR \& PARAMETER
BINDING}\label{walrus-operator-parameter-binding}}

\hypertarget{the-assignment-expression-pep-572}{%
\subsection{\texorpdfstring{14.1 The Assignment Expression (\texttt{:=}
/ PEP
572)}{14.1 The Assignment Expression (:= / PEP 572)}}\label{the-assignment-expression-pep-572}}

Introduced in Python 3.8, the assignment expression operator
(colloquially called the \textbf{walrus operator}) allows assigning
values to variables inside expressions. This departs from Python's
historical strict division between statements and expressions.

\hypertarget{ast-representation}{%
\subsubsection{1. AST Representation}\label{ast-representation}}

A standard assignment is a statement represented by the
\passthrough{\lstinline!Assign!} node in the Abstract Syntax Tree (AST),
which does not return a value. The walrus operator compiles to a
\passthrough{\lstinline!NamedExpr!} node, which is an expression node
returning the bound value.

\begin{lstlisting}
      [Assign Statement]                  [NamedExpr Expression]
         /         \                            /         \
   Name(x)      Constant(1)               Name(x)      Constant(1)
   (Returns void/no value)              (Returns 1 to parent context)
\end{lstlisting}

\hypertarget{bytecode-stack-operations}{%
\subsubsection{2. Bytecode \& Stack
Operations}\label{bytecode-stack-operations}}

Let's examine how CPython processes a normal assignment vs.~an
assignment expression. Consider the following source codes and their
corresponding bytecodes:

\begin{lstlisting}[language=Python]
# Standard Assignment
x = 1
\end{lstlisting}

\begin{lstlisting}
2 LOAD_CONST               0 (1)
4 STORE_FAST               0 (x)
\end{lstlisting}

\begin{lstlisting}[language=Python]
# Assignment Expression
(x := 1)
\end{lstlisting}

\begin{lstlisting}
2 LOAD_CONST               0 (1)
4 DUP_TOP
6 STORE_FAST               0 (x)
\end{lstlisting}

\emph{Note: In Python 3.11+, the compiler utilizes specialized stack
manipulation instructions (like \passthrough{\lstinline!COPY 1!})
instead of \passthrough{\lstinline!DUP\_TOP!}, but the mechanism remains
identical.}

The execution steps for the walrus operator stack operation are: 1.
\passthrough{\lstinline!LOAD\_CONST!} pushes the reference to the
constant integer \passthrough{\lstinline!1!} onto the value stack. 2.
\passthrough{\lstinline!DUP\_TOP!} replicates the top element of the
stack, resulting in two references to \passthrough{\lstinline!1!} on the
stack. 3. \passthrough{\lstinline!STORE\_FAST!} pops one of the
references and binds it to the local variable name
\passthrough{\lstinline!x!}. 4. The remaining reference to
\passthrough{\lstinline!1!} remains on the top of the stack, allowing
parent expressions (such as \passthrough{\lstinline!if!} conditionals or
loops) to consume it.

\hypertarget{scoping-rules-and-symbol-table-traversal}{%
\subsubsection{3. Scoping Rules and Symbol Table
Traversal}\label{scoping-rules-and-symbol-table-traversal}}

To prevent variable leakage and namespace pollution, PEP 572 specifies
complex scoping rules, especially within comprehensions: *
\textbf{Comprehensions}: Python list, set, and dict comprehensions, as
well as generator expressions, are executed in a nested function scope.
* \textbf{Variable Binding}: An assignment expression
\passthrough{\lstinline!x := ...!} inside a comprehension binds the
variable \passthrough{\lstinline!x!} in the \emph{surrounding} (parent)
scope, not the local comprehension scope. * \textbf{Symbol Table
Resolution}: During compilation, the symbol table builder
(\passthrough{\lstinline!symtable.c!}) traverses variables inside
comprehensions. When it encounters a
\passthrough{\lstinline!NamedExpr!}, it marks the target variable as
free/cell or nonlocal/global depending on the parent scope, hoisting the
reference out of the comprehension's inner namespace. *
\textbf{Restrictions}: A target variable cannot share a name with a
\passthrough{\lstinline!.0!} parameter (the iterator variable) or an
explicit comprehension loop target (e.g.,
\passthrough{\lstinline![i for i in range(10) if (i := 2)]!} is a syntax
error).

\begin{center}\rule{0.5\linewidth}{0.5pt}\end{center}

\hypertarget{positional-only-parameters-pep-570}{%
\subsection{\texorpdfstring{14.2 Positional-Only Parameters (\texttt{/}
/ PEP
570)}{14.2 Positional-Only Parameters (/ / PEP 570)}}\label{positional-only-parameters-pep-570}}

Python 3.8 introduced the \passthrough{\lstinline!/!} marker to define
positional-only parameters. Any parameters declared before the
\passthrough{\lstinline!/!} marker cannot be passed as keyword
arguments.

\hypertarget{c-level-representation-in-pycodeobject}{%
\subsubsection{\texorpdfstring{1. C-level Representation in
\texttt{PyCodeObject}}{1. C-level Representation in PyCodeObject}}\label{c-level-representation-in-pycodeobject}}

The compiler encodes parameter boundaries directly into the function's
bytecode descriptor struct. In \passthrough{\lstinline!code.h!}, the
\passthrough{\lstinline!PyCodeObject!} structure contains specific
integer fields:

\begin{lstlisting}[language=C]
typedef struct {
    PyObject_HEAD
    int co_argcount;            /* Number of positional arguments */
    int co_posonlyargcount;     /* Number of positional-only arguments */
    int co_kwonlyargcount;      /* Number of keyword-only arguments */
    /* ... additional fields ... */
} PyCodeObject;
\end{lstlisting}

\hypertarget{argument-parsing-calling-execution-path}{%
\subsubsection{2. Argument Parsing \& Calling Execution
Path}\label{argument-parsing-calling-execution-path}}

When a function is called, CPython passes arguments through the internal
calling mechanism (\passthrough{\lstinline!\_PyEval\_EvalCodeWithName!}
or \passthrough{\lstinline!\_PyFunction\_Vectorcall!}). 1. If the caller
provides arguments via keyword dict mappings, the interpreter extracts
the argument name and performs lookup check. 2. The interpreter compares
the number of passed positional arguments against
\passthrough{\lstinline!co\_posonlyargcount!}. 3. If keyword arguments
matching a parameter name before
\passthrough{\lstinline!co\_posonlyargcount!} are detected, the
interpreter raises a runtime \passthrough{\lstinline!TypeError!}:
\passthrough{\lstinline!TypeError: f() got some positional-only arguments passed as keyword arguments!}.

\hypertarget{micro-performance-optimization}{%
\subsubsection{3. Micro-Performance
Optimization}\label{micro-performance-optimization}}

Using positional-only parameters offers clear performance gains: 1.
\textbf{Bypassing Hash Lookups}: When keyword arguments are eliminated,
the interpreter bypasses dictionary keys extraction and hash checks for
the positional-only segment of the arguments tuple. 2. \textbf{Direct
Array Writes}: The VM copies references directly from the incoming
contiguous array of positional arguments (vectorcall array) straight
into the local frame variable storage.
\[\text{Offset Index} = \text{Passed Index}\]

\begin{center}\rule{0.5\linewidth}{0.5pt}\end{center}
